\chapter{System Requirements and Threat Model}\label{ch:req}

\section{Functional requirements}

The requirements below are those the firmware actually implements; each is traceable to the
source file and function given in Chapter~\ref{ch:impl}.

\begin{longtable}{@{}p{0.07\textwidth}p{0.60\textwidth}p{0.26\textwidth}@{}}
\toprule
\textbf{ID} & \textbf{Requirement} & \textbf{Where implemented} \\
\midrule
\endhead
FR1 & Detect motion through a passive infrared sensor, with a warm-up period and a
confirmation window to suppress spurious transitions & \code{readSensors()} \\
FR2 & Detect door or window opening through a switch, debounced & \code{readSensors()} \\
FR3 & Arm and disarm locally through a momentary button, debounced & \code{readSensors()},
\code{stepFsm()} \\
FR4 & Maintain a five-state alarm state machine: disarmed, exit delay, armed, entry delay,
alarm & \code{stepFsm()} \\
FR5 & Drive a siren, two status indicators and a local display on state change &
\code{driveOutputs()} \\
FR6 & Retain the armed state across a power interruption & NVS, \code{persistState()} \\
FR7 & Accept remote commands \code{ARM}, \code{ARM\_AWAY}, \code{DISARM} and \code{STATUS}
over MQTT & \code{handleCommands()} \\
FR8 & Authenticate remote commands with a message authentication code over a monotonic counter
& \code{authmac.h}, \code{handleCommands()} \\
FR9 & Publish state and events, with a retained state topic and a periodic heartbeat &
\code{networkTask()} \\
FR10 & Announce unavailability through an MQTT Last Will message & \code{networkTask()} \\
FR11 & Continue to perform FR1--FR6 with no network connectivity & core-1 loop, by
construction \\
\bottomrule
\end{longtable}

FR11 is the requirement that shapes the architecture. It is the reason the safety function
cannot be permitted to depend on, wait for, or be scheduled behind any network operation.

\section{Non-functional requirements}

\begin{longtable}{@{}p{0.07\textwidth}p{0.42\textwidth}p{0.44\textwidth}@{}}
\toprule
\textbf{ID} & \textbf{Requirement} & \textbf{Status in this work} \\
\midrule
\endhead
NFR1 & Sensor-to-siren latency bounded by an application deadline & bound \emph{derived}
analytically (Section~\ref{sec:impl-rt}); \textbf{not measured} \\
NFR2 & The safety loop never blocks indefinitely & established: zero-timeout mutex acquisition
and a bounded drain, after the repair in Section~\ref{sec:impl-drain} \\
NFR3 & Network load must not alter the latency distribution & \emph{logical} isolation
established by model checking; \textbf{timing isolation not measured} \\
NFR4 & Remote commands must resist capture and replay & established: monotonic counter
persisted before the command is acted on \\
NFR5 & Legitimate operator commands must survive an attack & measured: 1,250 of 1,250 admitted
under flood with authentication \\
NFR6 & The firmware must fit the device's flash and RAM & measured at build time
(Section~\ref{sec:impl-cost}) \\
NFR7 & All experimental results must be reproducible from committed artefacts & established:
byte-identical regeneration from fixed seeds \\
\bottomrule
\end{longtable}

NFR1 and NFR3 are the two requirements this work does \emph{not} discharge, and the distinction
between them and the rest is maintained throughout the report.

\section{System model}

The device is a dual-core microcontroller with one integrated radio. Write $C_{s}$ for the core
hosting the safety function and $C_{n}$ for the core hosting all network activity. The safety
function is a periodic loop that samples sensors, advances a finite state machine and drives
outputs.

Let $t_{d}$ be the instant at which a physical intrusion asserts a sensor line, and $t_{a}$ the
instant at which the siren is driven. The quantity of interest is the sensor-to-siren latency
\begin{equation}
L = t_{a} - t_{d},
\label{eq:rep-latency}
\end{equation}
and the safety requirement is that $L \le L_{\max}$ for a deadline $L_{\max}$ fixed by the
application.

Let $\Lambda$ denote the offered load on the radio in frames per second, and $F_{\Lambda}$ the
distribution of $L$ under that load. The isolation property one wishes to establish is
\begin{equation}
F_{\Lambda} \stackrel{d}{=} F_{0} \quad \text{for all } \Lambda \text{ in the operating range},
\label{eq:rep-invariance}
\end{equation}
that is, that the latency distribution under load is indistinguishable from the unloaded
control.

\subsection{Two claims requiring different evidence}

Equation~\eqref{eq:rep-invariance} decomposes into two claims, and keeping them apart is the
single most important methodological decision in this work.

\begin{description}[leftmargin=1.8em,itemsep=4pt]
  \item[Logical isolation.] No control- or data-flow dependency exists by which network
        activity can block, delay or corrupt the safety path. This is a property of the
        \emph{program}, and it is decidable by model checking and simulation.
  \item[Timing isolation.] Contention for shared silicon — the memory bus, caches, the
        interrupt controller — does not shift $F_{\Lambda}$ beyond the deadline margin. This is
        a property of the \emph{implementation on particular hardware}, and it is decidable only
        by measurement.
\end{description}

Logical isolation is necessary but not sufficient for \eqref{eq:rep-invariance}. This project
establishes the former and specifies, but does not execute, the experiment for the latter. Every
claim in this report about isolation is a claim about logical isolation unless it explicitly
states otherwise.

\section{Security objectives}

\begin{enumerate}[leftmargin=2.4em,itemsep=4pt,label=\textbf{SO\arabic*.}]
  \item An adversary without credentials must not be able to change the alarm's state.
  \item An adversary who observes a legitimate command must not be able to replay it.
  \item An adversary must not be able to prevent a legitimate operator command from taking
        effect.
  \item An adversary must not be able to suppress the alarm's response to a physical intrusion.
  \item A build configuration must not be able to enable an unauthenticated state change
        silently.
\end{enumerate}

SO5 is unusual as a stated objective and it is stated because this project failed it. The
defect, its measured consequence and its repair are the subject of
Section~\ref{sec:sec-guard}.

\section{Threat model}

\subsection{The primary adversary}

The primary adversary, designated \textbf{A}, is remote and has radio-layer access only: within
radio range of the device, able to transmit arbitrary 802.11 frames and to publish to the broker
as an unauthenticated client. Adversary A has \emph{no} credentials, no physical access and no
prior knowledge of any secret.

Unless a claim in this report explicitly names another scenario, it is a claim against
adversary A.

\subsection{Extended compromise scenarios}

Five further scenarios are evaluated separately, as extended compromise analyses. They are
reported in their own right and are never used to support a claim about the design.

\begin{table}[htbp]
\centering
\caption[Adversary scenarios]{Adversary scenarios. A is the primary threat model; B through
E$'$ are extended compromise analyses. The \emph{Access} column is load-bearing: E is the only
scenario requiring physical access to the device, and it is the one the layered configuration
survives. E$'$ requires none, and it is the one it does not.}
\label{tab:rep-threats}
\begin{tabular}{@{}llll@{}}
\toprule
& Additional capability & Access & Evaluated in \\
\midrule
A    & none (radio range only)             & remote            & throughout \\
B    & one shared broker credential        & remote            & Ch.~\ref{ch:security} \\
C    & controller credential               & remote            & Ch.~\ref{ch:security} \\
D    & device credential                   & remote            & Ch.~\ref{ch:security} \\
E    & device credential + command key     & \textbf{physical} & Ch.~\ref{ch:security} \\
E$'$ & controller credential + command key & \textbf{remote}   & Ch.~\ref{ch:security} \\
\bottomrule
\end{tabular}
\end{table}

\subsection{E and E$'$: why the distinction decides the result}

Two of the five are easily conflated, and the project's central access-control result turns on
the difference between them.

\textbf{Scenario E is a device flash dump.} It yields the device's broker credential and the
command key together, because both reside in the same flash image. Obtaining it requires
\emph{physical} access to the hardware — the adversary must hold the device.

\textbf{Scenario E$'$ is a compromised controller host} — the phone, the hub or the server that
legitimately issues commands. It yields a credential that the access control list \emph{must}
permit, together with the same command key, and it requires no physical access at all.

The layered configuration withstands E and falls to E$'$. The scenario that is cheaper for an
adversary to reach is therefore the one that defeats it. This is reported in
Section~\ref{sec:sec-eprime} and was found by measuring an arm that an earlier draft of the
work had not run; that earlier draft claimed the layered configuration ``survives every case'',
and the claim was wrong.

\subsection{Trust boundaries}

Four boundaries are relevant, and the architecture in Chapter~\ref{ch:arch} draws all four.

\begin{enumerate}[leftmargin=1.8em,itemsep=4pt]
  \item \textbf{Between the two cores.} The safety core trusts nothing the network core
        produces beyond a well-formed, authenticated command delivered through a bounded queue.
  \item \textbf{Between the device and the broker.} The device trusts the broker to deliver
        messages and \emph{not} to attest their origin. This is the boundary the configuration
        defect misread.
  \item \textbf{Between the broker and publishing clients.} The broker's access control list is
        the only mechanism distinguishing one authenticated client from another.
  \item \textbf{Between the device and its own flash.} An adversary with physical access
        crosses this boundary, which is what makes scenario E a physical-access scenario.
\end{enumerate}

\subsection{Authentication and authorisation assumptions}

\begin{itemize}[leftmargin=1.6em,itemsep=4pt]
  \item \textbf{Transport.} TLS is assumed to protect each hop against observation and
        modification. It is \emph{not} assumed to convey provenance across the broker, because
        the broker re-originates every message.
  \item \textbf{Broker authentication.} Clients present credentials from a password file.
        Possession of a valid credential establishes identity to the broker and nothing more.
  \item \textbf{Broker authorisation.} A deny-by-default access control list grants each
        identity specific publish and subscribe permissions per topic. An identity that must
        publish commands cannot be prevented from publishing commands.
  \item \textbf{Application-layer authentication.} A command is acted on only if its
        HMAC-SHA256 tag verifies under a shared secret and its counter exceeds the value
        persisted in non-volatile storage.
  \item \textbf{Key distribution is out of scope.} How the command key reaches the device and
        the controller is not addressed. This is a genuine limitation, not an oversight.
\end{itemize}

\subsection{Limitations of the threat model}

Four limitations should be read alongside every security claim in this report.

\begin{enumerate}[leftmargin=1.8em,itemsep=4pt]
  \item \textbf{No side-channel or fault-injection adversary is modelled.} Power analysis,
        electromagnetic analysis, glitching and similar techniques are entirely outside the
        model.
  \item \textbf{The broker itself is assumed honest.} A malicious or compromised broker can
        drop, delay or reorder messages. Message authentication prevents it forging a valid
        command, but availability is not protected against it.
  \item \textbf{Supply-chain compromise is not modelled.} The firmware, the toolchain and the
        libraries are assumed to be what they claim to be.
  \item \textbf{Denial of service at the radio layer is acknowledged but unmeasured.} The one
        adversary capability the model takes most seriously is the one whose consequence this
        work could not measure. Chapter~\ref{ch:limitations} does not minimise this.
\end{enumerate}
